\documentclass[11pt,a4paper]{article}
\usepackage[T1]{fontenc}
\usepackage[utf8]{inputenc}
\usepackage[english]{babel}
\usepackage{amsmath,amssymb}
\usepackage[margin=2cm]{geometry}
\setlength{\parindent}{0pt}
\setlength{\emergencystretch}{3em}
\allowdisplaybreaks

\title{About \texttt{document\_statements.tex}}
\author{}
\date{}
\begin{document}
\maketitle
This note explains what the test data in \texttt{document\_statements.tex} is used for. That file is input for the engine in \texttt{main.c}; it is read by the test \texttt{test\_output.c} in this folder.

\begin{enumerate}
\item \textbf{Purpose.} The file holds 5 statements, one of each kind the final document written by \texttt{output.c} has to report. The test is about the document, not about the mathematics: the statements are easy on purpose.

\item \textbf{What it exercises.} The writing of an item: the statement as it was read, then the proof or the solutions, or the verdict when there is no proof; and the graph, which becomes a page of the PDF of the graphs and is included in the document by page number.

\item \textbf{Variety.} A proof by induction, which has no graph; a solution set and a statement decided by a sign table, which have one each; a false statement, whose graph still shows where the condition holds; a statement that cannot be read, which has neither proof nor graph.

\item \textbf{Expected result.} The line \texttt{PASS: 25 of 25 output checks passed}.

\item \textbf{How to run it.} From the top folder of the repository:
\begin{verbatim}
cc -std=c11 -O2 -Wall -Wextra -o /tmp/test_output \
    tests/output/test_output.c && /tmp/test_output
\end{verbatim}
\end{enumerate}
\end{document}
